\section{Fast stability charts}\label{sec:charts}

\subsection{Brute-force sweep with interpolable coloring}\label{sec:coloring}

The single-point solve of Section~\ref{sec:method} is embarrassingly
parallel over the parameter plane; a multithreaded sweep over a grid
$\vect p_{ij}$ returns the fields $\Zint$, $\Zraw$, $\min|\Dfun|$, $\sest$
and $\west$.
The integer field $\Zint$ alone is a poor plotting quantity -- flat plateaus
separated by jumps -- but the estimated root location repairs precisely this
deficiency. The combined field
\begin{equation}\label{eq:coloring}
  C(\vect p) = \Zint(\vect p) + \bigl[\Zint(\vect p)=0\bigr]\cdot\sest(\vect p)
\end{equation}
-- where $[\cdot]$ denotes the Iverson bracket -- is a remarkably efficient
coloring trick for visual inspection: inside the
stable domain ($\Zint=0$, $\sest<0$) it varies \emph{smoothly} and displays
the spectral gap, i.e., how far the dominant root lies from the imaginary
axis, while outside it shows the integer number of unstable roots as
discrete color levels. One glance at a single map thus reveals where the
system is stable, how robustly stable it is there, and how many roots must
be stabilized elsewhere. Because $C$ is smooth over the stable domain,
interpolation between grid points is meaningful exactly where it is useful,
so a coarse and cheap grid suffices for the colored background; only the
integer levels of the unstable side are plotted as flat colors.

\subsection{High-resolution boundary tracing with MDBM}\label{sec:mdbm}

The stability boundary itself deserves better than pixel resolution. The
multi-dimensional bisection method (MDBM)~\cite{bachrathy2012bisection,bachrathy2025mdbm}
locates the zero set of a scalar function over a coarse initial mesh and
refines only the bracketing cells, halving the mesh size in every iteration.
Here \emph{only the estimated root distance is used} as the scalar objective:
\begin{equation}\label{eq:objective}
  g(\vect p) = \operatorname{sign}\!\bigl(\bigl[\Zint(\vect p)=0\bigr]-\tfrac12\bigr)\,
               \bigl|\sest(\vect p)-\sigma\bigr|,
\end{equation}
which is positive inside the stable ($\gamma$-stable) domain, negative
outside, crosses zero exactly on the stability boundary, and -- unlike
$\Zint$ itself -- varies almost linearly across it, because near the
boundary $\sest$ is an accurate estimate of the dominant root's real part
(Section~\ref{sec:tracking}). MDBM can therefore interpolate the zero
crossing within each bracketing cell to sub-cell accuracy. Note the contrast
with applying MDBM to the pair
$\Real\Dfun=\Imag\Dfun=0$ in the parameter--frequency space: that approach
traces \emph{all} root-crossing D-curves, most of which lie deep in the
unstable region (Section~\ref{sec:intro}), while the
objective~\eqref{eq:objective} follows only the \emph{true} stability limit.
By the very nature of the bisection refinement, the computational effort
concentrates exclusively on the boundary: $k$ refinement levels raise the
equivalent resolution of an $n_0\times n_0$ initial mesh to
$2^{k}n_0\times 2^{k}n_0$, while the number of function evaluations grows
only in proportion to the boundary length. On the fourth-order benchmark of
Section~\ref{sec:performance}, four refinement levels from a $20\times20$
mesh trace the boundary at a $320\times320$-equivalent resolution
(Table~\ref{tab:mdbm}) -- roughly two orders of magnitude fewer function
evaluations than a brute-force grid of that density.

\subsection{The most robust parameter point}\label{sec:robust}

With a high-resolution boundary and a smooth interior metric available, the
natural design question -- \emph{which parameter point is the most robustly
stable?} -- reduces to a geometric one: the center of the largest circle
inscribed in the stable domain maximizes the distance to the stability
boundary in parameter space. The package computes it from the MDBM boundary
point cloud (or from the binary stable mask of the brute-force grid) by a
coarse-to-fine search followed by an omnidirectional pattern refinement with
adaptive step control. Anisotropic parameter scaling turns the circle into
an ellipse, which avoids degenerate results in elongated stable regions and
lets the user weight the expected perturbations of the individual
parameters. The exact radius is rarely of engineering interest; the
\emph{location} of the optimum is, and that location is insensitive to the
residual discretization of the boundary. If a spectral-gap-optimal (rather
than parameter-robust) point is preferred, the interpolable interior field
$\sest(\vect p)$ of~\eqref{eq:coloring} can be minimized instead -- both
options are cheap post-processing steps on data already computed.

\subsection{$\sigma$-contours: decay-rate maps}\label{sec:sigmacontours}

Repeating the boundary tracing of Section~\ref{sec:mdbm} for a family of
shifted lines $\sigma\in\{0,\sigma_1,\sigma_2,\dots\}$ produces the level
curves of the guaranteed decay rate, i.e., a $\gamma$-stability map: each
contour bounds the parameter region whose rightmost root lies left of
$\Real\lam=\sigma$. The interior coloring by $\sest$ and the traced
$\sigma$-contours are two views of the same robustness landscape -- the
former is a fast estimate on a coarse grid, the latter a high-resolution
guaranteed boundary -- and plotting them together provides a strong visual
cross-validation of both (Section~\ref{sec:showcase}).

\subsection{Recommended workflow}\label{sec:workflow}

The complete chart-generation recipe used throughout the paper is:
\begin{enumerate}
  \item \emph{Coarse colored sweep.} Multithreaded evaluation
    of~\eqref{eq:phaseode} on a coarse grid; plot the combined
    field~\eqref{eq:coloring}. Cost: a fraction of a second to a few seconds.
  \item \emph{Boundary tracing.} MDBM on the scalar
    objective~\eqref{eq:objective}, typically $4$ refinement levels from a
    $20\times20$--$30\times30$ initial mesh. Cost: comparable to or below
    the sweep.
  \item \emph{Robust point.} Largest inscribed circle/ellipse from the MDBM
    point cloud. Cost: milliseconds.
  \item \emph{Optional.} $\sigma$-contour family; root refinement at
    selected points; higher tolerance for publication-quality boundaries.
\end{enumerate}
Every step reuses the same single-point solver, so the entire chart refreshes
fast enough for interactive parameter exploration
(Section~\ref{sec:performance}).